% One-time imported draft. Editable native source; no Markdown regeneration.
\section{Synchronous updates need a stable
  source}\label{008-life-generations-p01-synchronous-updates-need-a-stable-source}

Synchronous means that every next-state decision uses the same old
generation. The row and column loops run one instruction at a time, but
their decisions must act as if they were based on one shared snapshot.
The physical order of destination stores must not change what later
cells read.

An in-place sweep breaks that condition. In the default pattern, the
cell at (3,2) is born above the horizontal line. If its new state is
written into the old board, a later cell at (2,3) can count that newborn
along with its old neighbor at (3,3). It then sees two neighbors and
survives, although the synchronous rule would see only one and kill it.
The local rule can be implemented perfectly and still receive the wrong
generation as input.

Separate current and destination allocations prevent that dependency.
Each result goes to the other board. Only after every cell is complete
do the two allocations exchange roles. The distinction will recur in
image filters, finite-difference stencils, and other computations whose
outputs must not contaminate their inputs.
